\section{Structured README File}
In this section, we present our first attempt to facilitate LLM-based code generation under domain specific requirements.
To improve the ability of LLMs to understand and correctly use the RDPro library, we propose redesigning the README file into a more structured and pattern consistent format. The central idea is to provide documentation that is explicitly organized and semantically regular, enabling LLMs to reliably infer usage patterns across functions.

In the original documentation, function descriptions are often surrounded by extensive notes, informal explanations, and loosely organized examples. While such documentation may be helpful for human readers, it introduces ambiguity and noise when LLM start to run code generation and reasoning. A more structured format reduces this ambiguity and provides clearer signals about functional intent and usage. The revised README begins with a background section that introduces the relevant domain context (e.g., raster data format), followed by a concise explanation of the core data model and metadata abstractions used in RDPro. For each callable function, the documentation should follow a consistent schema that includes: 
(1) \textbf{Function Name};
(2) \textbf{Description} (high level purpose and semantics);
(3) \textbf{Input Parameters} (data types, constraints, and meanings);
(4) \textbf{Output} (return type and semantics);
(5) \textbf{Example Usage} (minimal, executable code snippet)
This structured template enables LLMs to learn reusable documentation patterns, map natural language requests to function signatures, and generate syntactically and semantically correct code with higher reliability.